\section{Particles and reservations}
\label{sec:particles}

We next construct consistent bag laws while preserving selected cut measures
along specified portions of the decomposition.  The basic operation changes
two functions without changing their combined trace on the current bag.  A
uniform bound on all functions survives arbitrarily many such operations.

Throughout this section, let $\Delta=\operatorname{diam}(V)>0$, and fix a
rooted tree decomposition whose bags have size at most $k$, where $k\ge2$.
A \emph{particle} is a real function on all of $V$.  We maintain a list
$F_1,\ldots,F_N$ of particles, initially all zero.  The sweep
$\sigma_F^M$ and symmetrization of a cut measure have the meanings given in
Section~\ref{sec:cuts}; in particular, $\sigma_F^M$ has mass $2M$.

If a common cutoff $M$ contains all current particle heights, the
measure carried by the list has trace
\[
 \left.\sum_{i=1}^N(\sigma_{F_i}^M)^{\mathrm{sym}}\right|_P
\]
on a bag $P$. The update below preserves this entire trace, including
the constant assignments. We will choose one cutoff that remains valid
through every update after proving uniform bounds on the heights.

\begin{lemma}[Affine trace update]
\label{lem:trace-update}
Let $P\subseteq V$, and suppose that particles $F_1,F_2$ admit a common
weak increasing order on $P$.  Put $f=(F_1+F_2)/2$, and let $g:V\to\mathbb R$
vanish on $P$.  If the four functions $F_1,F_2,f+g,f-g$ take values in
$[-M,M]$, then
\[
 \left.(\sigma_{F_1}^M+\sigma_{F_2}^M)\right|_P
 =\left.(\sigma_{f+g}^M+\sigma_{f-g}^M)\right|_P.
\]
The equality remains true after symmetrization.  For nonempty $P$, it
includes the masses of both constant assignments on $P$.
\end{lemma}

\begin{proof}
Write $P=\{p_1,\ldots,p_r\}$ in a common weak increasing order.  For a
function $F$ with that order, sweeping the labels
$\mathbf 1_{\{F(p)>t\}}$ gives, successively, the monotone assignments with
masses
\[
 F(p_1)+M,\qquad
 F(p_{i+1})-F(p_i)\quad(1\le i<r),\qquad
 M-F(p_r).
\]
Every other assignment has mass zero.  These expressions are affine in the
height vector, including the first and last expressions for the constant
assignments.  Equal heights merely give zero masses.  Thus the sum of the
traces for $F_1$ and $F_2$ is twice the trace for $f$.  Since $f+g$ and
$f-g$ both equal $f$ on $P$, the assertion follows.  For $P=\varnothing$
there is just one assignment, of mass $2M$ per sweep, and the same
conclusion holds.  Symmetrization and restriction are linear operations,
so they preserve the identity.
\end{proof}

An \emph{update at $P$} replaces two particles having a common weak order
on $P$ by
\[
 F_i,F_j\ \longmapsto\ f+g,f-g,
 \qquad f=\frac{F_i+F_j}{2},\qquad g|_P=0.
\]
We call $g$ the perturbation.  The empty set is an allowed interface before
the root bag is processed.  Fix constants $L,B>0$ and require every
perturbation to satisfy
\[
 \operatorname{Lip}(g)\le L,
 \qquad
 \sup_{v\in V}|g(v)|\le B\Delta.
\]
The choices $L=B=3$ will suffice in all our applications.  Keeping the two
parameters separate makes clear which bound is used at each step.

Applying the triangle inequality after each update would increase the
maximum Lipschitz bound by $L$ at every step.  Instead, for each fixed
$u\ne v$, we control all the normalized differences
$(F_i(u)-F_i(v))/(L d(u,v))$ together.  The sum over any $r$ indices is
bounded by $r(N-r)$.  Averaging the bounds for $r-1$ and $r+1$ leaves
exactly one unit for the new perturbation, so the bound does not
increase with the number of updates.

\begin{lemma}[Uniform particle bounds]
\label{lem:particle-bounds}
After any finite sequence of updates starting from zero particles,
\[
 \operatorname{Lip}(F_i)\le L(N-1),
 \qquad
 |F_i(v)|\le B(N-1)\Delta
 \quad (1\le i\le N,\ v\in V).
\]
These bounds hold for arbitrary choices of the updated pairs and
perturbations subject to the displayed bounds, regardless of the number
of updates.
\end{lemma}

\begin{proof}
We first prove a numerical invariant.  Suppose a vector
$a=(a_1,\ldots,a_N)$ starts at zero and is changed by operations of the
form
\[
 a_i,a_j\ \longmapsto\frac{a_i+a_j}{2}+e,
                         \frac{a_i+a_j}{2}-e,
 \qquad |e|\le1.
\]
At every stage its total is zero, and
\[
 \sum_{i\in I}a_i\le |I|(N-|I|)
 \qquad\text{for every }I\subseteq\{1,\ldots,N\}.
\]
For the induction, put $Q_r=r(N-r)$, including $Q_0=Q_N=0$.
The assertion is true initially.  An update leaves the sum over $I$
unchanged when $I$ contains both updated indices or neither.  If $I$
contains exactly one of them, write $S$ for the remaining $r-1$ indices,
where $r=|I|$.  Its new sum is at most
\[
 \frac12\left(\sum_{h\in S}a_h+
                   \sum_{h\in S\cup\{i,j\}}a_h\right)+1
 \le\frac12(Q_{r-1}+Q_{r+1})+1
 =Q_r.
\]
This proves the invariant.  Singleton subsets give $a_i\le N-1$; their
complements, together with the zero total, give $a_i\ge-(N-1)$.

For distinct $u,v$, apply the invariant to
\[
 a_i=\frac{F_i(u)-F_i(v)}{L\,d(u,v)}.
\]
The corresponding perturbation is
$e=(g(u)-g(v))/(L\,d(u,v))$, of absolute value at most one.  This proves
the Lipschitz bound.  For each fixed $v$, apply the same invariant to
$a_i=F_i(v)/(B\Delta)$, with $e=g(v)/(B\Delta)$, to obtain the value
bound.  The argument applies separately to every branch of a copied
particle history.
\end{proof}

A \emph{reservation} freezes a pair of particles: neither particle may
be used in another update while that reservation is in force.  Releasing
a reservation makes its particles available again; it does not remove
them from the list.  A reservation may be continued into any prescribed
children using copies of the same pair.  At a branching bag, the whole
list and all reservation records are copied separately for each child.
Consequently the number of simultaneous reservations is counted within
one copy of the memory, not across different branches.

The following proposition packages the construction.  It allows an
arbitrary finite schedule of launches and releases; applications will
choose their schedules from the graph and decomposition alone.
The parameter $q$ counts only pairs held at the same time in one copy
of the list.  This is the quantity that must be fixed before choosing
any geometric scale parameters.

\begin{proposition}[Realization by particles]
\label{prop:particle-realization}
Suppose a finite schedule requires at most $q$ simultaneous reserved
pairs in each memory copy, including any temporary reservation made
before another is released.  Updates may be made at the current bag, or
on a child-specific copy at its parent bag before moving to the child.
Before processing the root, updates at the empty interface are also
allowed.  Every update uses a perturbation satisfying the bounds above
and vanishing on its interface.

Choose
\[
 N=2q+k!+2,
 \qquad
 M=\bigl(B(N-1)+2\bigr)\Delta.
\]
Then all requested updates and reservations can be carried out.  They
produce symmetric bag measures $\mu_t$ with the following properties.
\begin{enumerate}
\item Each $\mu_t$ has the same total mass $m=2NM$, and neighboring bag
measures have equal traces on their intersection.
\item Every particle height is at distance at least $2\Delta$ from each
endpoint of the sweep interval $[-M,M]$.
\item The pseudometric $D_\mu$ obtained by Markov gluing satisfies
\[
 D_\mu(u,v)\le U\,d(u,v)
 \quad\text{for all }u,v\in V,
 \qquad U=NL(N-1).
\]
\item If a pair is reserved from a launch bag through a later bag on a
descending path, its fixed symmetric two-particle cut measure has trace
dominated by $\mu_t$ at every bag on that segment.  The segment includes
the launch and release bags, even when the particles are reused at the
release bag.  A reservation launched at the empty interface has the same
property from the root onward.
\end{enumerate}
The statements hold for each realization of any auxiliary randomness.
For measurable randomized choices, taking expectations also gives an
$L_1$ pseudometric with the same upper bound.
\end{proposition}

\begin{proof}
At every stage at most $2q$ particles are reserved, leaving at least
$k!+2$ free particles.  For a bag of size $r\le k$, choose a weak
increasing listing for each free particle, breaking ties by a fixed
ordering of the vertices.  There are at most $r!\le k!$ such listings,
so two free particles have the same listing.  They are eligible for the
next update.  At the empty interface every pair is eligible.  This
also permits updates that do not initiate a reservation.
Fix an ordering of particle pairs and always choose the first eligible
pair using the pre-launch memory, before sampling any fresh perturbation
randomness for that launch.

Lemma~\ref{lem:particle-bounds} applies to each branch history, including
the operations performed on copies at the parent bag.  In particular,
the stated choice of $M$ gives the asserted $2\Delta$ margins at every
stage.  It is fixed before the construction begins and therefore is the
same for every particle, bag, and branch.

Whenever the current bag is $B_t$, let its recorded law be
\[
 \mu_t=
 \left.\sum_{i=1}^N(\sigma_{F_i}^M)^{\mathrm{sym}}\right|_{B_t}.
\]
Updates at $B_t$ leave this measure unchanged by
Lemma~\ref{lem:trace-update}, and releasing reservations makes no change
to the functions.  Thus it is immaterial whether this law is recorded
before or after any operations at that bag.  Preliminary operations at
the empty interface take place before the root law is recorded.

For a child $t'$ of $t$, copy the current memory and perform any
child-specific operations while still at $B_t$.  Its trace on $B_t$
remains $\mu_t$.  Move that copy to $B_{t'}$ and record its law there.
The two laws agree on $B_t\cap B_{t'}$, since both are restrictions of
the same total cut measure carried by this copy before its operations
at $B_{t'}$.  Those subsequent operations preserve the child law.
This proves consistency on every decomposition edge.  Each summand has
mass $2M$, both restriction and symmetrization preserve mass, and hence
every bag law has mass $2NM$.  The construction of
Section~\ref{sec:cuts} therefore supplies the glued cut pseudometric $D_\mu$.

If $uv$ is a graph edge, choose a bag containing both endpoints.  The
glued law has the prescribed marginal on that bag.  Sweeping separates
$u$ and $v$ with mass $|F_i(u)-F_i(v)|$, and symmetrization preserves
that mass.  Consequently, by Lemma~\ref{lem:particle-bounds},
\[
 D_\mu(u,v)
 =\sum_{i=1}^N|F_i(u)-F_i(v)|
 \le NL(N-1)d(u,v)
 \le U\ell(uv).
\]
Applying the triangle inequality along a shortest graph path proves
$D_\mu(u,v)\le U d(u,v)$ for arbitrary vertices.

For the last assertion, denote the two functions at launch by $H_1,H_2$
and set
\[
 \nu=(\sigma_{H_1}^M)^{\mathrm{sym}}
       +(\sigma_{H_2}^M)^{\mathrm{sym}}.
\]
They remain unchanged while reserved, including in each child copy
where the reservation continues.  At any bag in that portion of the
path, their two traces are summands of the current total trace, so
$\nu|_{B_t}\le\mu_t$.  At the launch bag this statement can be read just
after the launching update; at the release bag it can be read just
before the release.  All other operations at either bag preserve the
total trace.  Thus the same domination holds for its recorded law even
if the pair is immediately reused there.  The case of a reservation
started before the root is identical from its first bag onward.

All bounds and trace identities hold for each realization of the
auxiliary randomness.  For measurable choices, expectation of the
resulting finite cut measures realizes $\mathbb E D_\mu$ as an $L_1$
pseudometric, as in Section~\ref{sec:cuts}, and preserves the displayed
upper inequality.
\end{proof}

The number of updates and the total number of reservations over the
whole tree do not occur in $U$.  With the common choice $L=B=3$, one may
use the explicit bounds
\[
 \operatorname{Lip}(F_i)\le3(N-1),\qquad
 M=(3N-1)\Delta,\qquad U=3N(N-1).
\]
These constants depend only on $k$ and the maximum simultaneous
reservation count $q$.
